% -*- coding: utf-8 -*-
% !TEX program = xelatex
\documentclass[plain,noanswer,medmath]{../../template/jnuexam}
\usepackage{../../template/kysx}

\begin{document}


\examtitle{1989}{4}


\loadproblems{1}{1989P1}%载入试卷一


\makepart{填空题}{本题满分15分，每小题3分}


\begin{problem}
曲线$y=x+\sin^2 x$在点$\left(\frac{\pi}{2},1+\frac{\pi}{2}\right)$处的切线方程是\fillin{}.
\end{problem}


\begin{problem}
幂级数$\sum_{n=0}^{\infty}\frac{x^n}{\sqrt{n+1}}$的收敛域是\fillin{}.
\end{problem}


\begin{problem}
齐次线性方程组$\begin{cases}
  \lambda x_1+x_2+x_3=0, \\
  x_1+\lambda x_2+x_3=0, \\
  x_1+x_2+x_3=0 \\
\end{cases}$只有零解，则$\lambda$应满足的条件是\fillin{}.
\end{problem}


\begin{problem}
设随机变量$X$的分布函数为
$$F(x)=\begin{cases}
        0, & x<0, \\
  A\sin x, & 0\le x\le \frac{\pi}{2}, \\
        1, & x>\frac{\pi}{2},
\end{cases}$$
则$A=$\fillin{}，$P\left\{|X|<\frac{\pi}{6}\right\}=$\fillin{}．
\end{problem}


\begin{problem}
设随机变量$X$的数学期望$E(X)=\mu$，方差$D(X)=\sigma^2$,
则由切比雪夫(Chebyshev)不等式，有$P\{|X-\mu|\ge 3\sigma\}\le$ \fillin{}．
\end{problem}


\makepart{选择题}{本题满分15分，每小题3分}

\begin{problem}
设$f(x)=2^x+3^x-2$，则当$x\to 0$时 \pickout{}
\begin{abcd}
\item $f(x)$与$x$是等价无穷小量．
\item $f(x)$与$x$是同阶但非等价无穷小量．
\item $f(x)$是比$x$较高阶的无穷小量．
\item $f(x)$是比$x$较低阶的无穷小量．
\end{abcd}
\end{problem}


\begin{problem}
在下列等式中，正确的结果是 \pickout{}
\begin{abcd}
\item $\int f'(x)\dx=f(x)$．
\item $\int \d f(x)=f(x)$．
\item $\frac{\d}{\dx}\int f(x)\dx=f(x)$．
\item $\d\int f(x)\dx=f(x)$．
\end{abcd}
\end{problem}


\useproblem{1}{2}{5}


\begin{problem}
设$A$和$B$均为$n\times n$矩阵，则必有 \pickout{}
\begin{abcd}
\item $|A+B|=|A|+|B|$．
\item $AB=BA$．
\item $|AB|=|BA|$．
\item $(A+B)^{-1}=A^{-1}+B^{-1}$．
\end{abcd}
\end{problem}


\begin{problem}
以$A$表示事件“甲种产品畅销，乙种产品滞销”，则其对立事件$\overline{A}$为 \pickout{}
\begin{abcd}
\item “甲种产品滞销，乙种产品畅销”．
\item “甲、乙两种产品均畅销”．
\item “甲种产品滞销”．
\item “甲种产品滞销或乙种产品畅销”．
\end{abcd}
\end{problem}


\makepart{计算题}{本题满分15分，每小题5分}

\begin{problem}
求极限$\lim_{x\to \infty}\left(\sin\frac{1}{x}+\cos\frac{1}{x}\right)^x$．
\end{problem}


\begin{problem}
已知$z=f(u,v),u=x+y,v=xy$，且$f(u,v)$的二阶偏导数都连续．
求$\frac{\pd^2 z}{\pdx\pd y}$．
\end{problem}


\begin{problem}
求微分方程$y''+5y'+6y=2\e^{-x}$的通解．
\end{problem}


\makepart{}{本题满分9分}

\begin{problem}
设某厂家打算生产一批商品投放市场．已知该商品的需求函数为
$$P=P(x)=10\e^{-\frac{x}{2}},$$
且最大需求量为$6$，其中$x$表示需求量，$P$表示价格．
\begin{enumerate}
\item 求该商品的收益函数和边际收益函数．%(2分)
\item 求使收益最大时的产量、最大收益和相应的价格．%(4分)
\item 画出收益函数的图形．%(3分)
\end{enumerate}
\end{problem}


\makepart{}{本题满分9分}

\begin{problem}
已知函数$f(x)=\begin{cases}
    x, & 0\le x\le 1, \\
  2-x, & 1\le x\le 2. \\
\end{cases}$ 试计算下列各题：\par
\begin{enumerate*}[itemjoin=\relax]
\item $S_0=\int_0^2f(x)\e^{-x}\dx$；\hspace*{1em}%(4分)
\item $S_1=\int_2^4f(x-2)\e^{-x}\dx$；\\%(2分)
\item $S_n=\int_{2n}^{2n+2}f(x-2n)\e^{-x}\dx(n=2,3,\cdots)$；\hspace*{1em}%(1分)
\item $S=\sum_{n=0}^{\infty}S_n$．%(2分)
\end{enumerate*}
\end{problem}


\makepart{}{本题满分6分}

\begin{problem}
假设函数$f(x)$在$[a,b]$上连续，在$(a,b)$内可导，且$f'(x)\le 0$，记
$$F(x)=\frac{1}{x-a}\int_a^x f(t)\dt,$$
证明在$(a,b)$内，$F'(x)\le 0$．
\end{problem}


\makepart{}{本题满分5分}

\begin{problem}
已知$X=AX+B$，其中$A=\begin{pmatrix}
   0 & 1 & 0 \\
  -1 & 1 & 1 \\
  -1 & 0 & -1
\end{pmatrix}$, $B=\begin{pmatrix}
  1 & -1 \\
  2 & 0 \\
  5 & -3
\end{pmatrix}$，求矩阵$X$．
\end{problem}


\makepart{}{本题满分6分}

\begin{problem}
设$\alpha_1=(1,1,1)$, $\alpha_2=(1,2,3)$, $\alpha_3=(1,3,t)$．
\begin{enumerate}
\item 问当$t$为何值时，向量组$\alpha_1,\alpha_2,\alpha_3$线性无关？%(3分)
\item 问当$t$为何值时，向量组$\alpha_1,\alpha_2,\alpha_3$线性相关？%(1分)
\item 当向量组$\alpha_1,\alpha_2,\alpha_3$线性相关时，
将$\alpha_3$表示为$\alpha_1$和$\alpha_2$的线性组合．%(2分)
\end{enumerate}
\end{problem}


\makepart{}{本题满分5分}

\begin{problem}
设$A=\begin{pmatrix}
  -1 &  2 & 2 \\
   2 & -1 & -2 \\
   2 & -2 & -1
\end{pmatrix}.$
\begin{enumerate}
\item 试求矩阵$A$的特征值；%(2分)
\item 利用(I)小题的结果，求矩阵$E+ A^{-1}$的特征值，其中$E$是$3$阶单位矩阵．%(3分)
\end{enumerate}
\end{problem}


\makepart{}{本题满分7分}

\begin{problem}
已知随机变量$X$和$Y$的联合密度为
$$f(x,y)=\begin{cases}
 \e^{-(x+y)}, & 0<x<+\infty, 0<y<+\infty, \\
           0, & \text{其他}. \\
\end{cases}$$
试求：
\begin{enumerate*}
\item $P\{X<Y\}$；%(5分)
\item $E(XY)$．%(2分)
\end{enumerate*}
\end{problem}


\makepart{}{本题满分8分}

\begin{problem}
设随机变量$X$在$[2,5]$上服从均匀分布，现在对$X$进行三次独立观测，
试求至少有两次观测值大于$3$的概率．
\end{problem}


\end{document}
